% GENERATED, edit the French unit or the translation table
\mxunit{expertes-c08}{Prime numbers and Fermat's little theorem}
\mxsubtitle{Infinitude of the primes, factorisation, Fermat's little theorem}

\begin{objectives}
Recognise a prime number; prove that there are infinitely many primes; factorise an integer into primes and use the factorisation (divisors, gcd); use Fermat's little theorem to compute remainders.
\end{objectives}
\begin{prerequisites}
Chapters 6 and 7; binomial coefficients; binomial formula (chapter 1).
\end{prerequisites}

\section{Fluency}
\begin{fluency}[No calculator, five minutes]
\begin{enumerate}
\item Is the integer $91$ prime?
\item Give the prime numbers less than $30$.
\item Factorise $360$ into primes.
\item Factorise $\num{2026}$ into primes, given that $\num{1013}$ is prime.
\item Compute $2^{10}$ and give its remainder in the division by $11$.
\item Are the integers $35$ and $72$ coprime?
\end{enumerate}
\end{fluency}
\begin{solution}[Fluency]
\begin{enumerate}
\item No: $91=7\times13$.
\item $2$, $3$, $5$, $7$, $11$, $13$, $17$, $19$, $23$, $29$.
\item $360=2^3\times3^2\times5$.
\item $\num{2026}=2\times\num{1013}$.
\item $2^{10}=\num{1024}=11\times93+1$: the remainder is $1$.
\item Yes: $35=5\times7$ and $72=2^3\times3^2$ have no prime factor in common.
\end{enumerate}
\end{solution}

\mxnewpage
\section{Introductory activity}
\begin{activity}[A pattern in powers \mxchip{35 min}]
For an integer $n\ge2$ and an integer $a$, let $r_n(a)$ be the remainder of $a^{n-1}$ in the division by $n$.
\begin{enumerate}
\item Compute $r_5(2)$, $r_5(3)$, $r_7(2)$, $r_7(3)$ by hand, using congruences.
\item Using the Python function \texttt{pow(a, n - 1, n)}, which returns this remainder, complete the following table.
\begin{center}\mxarraystretch{1.8}
\begin{tabular}{c|*{7}{p{9mm}}}
$n$ & $5$ & $6$ & $7$ & $8$ & $9$ & $10$ & $11$\\\hline
$r_n(2)$ & & & & & & & \\
$r_n(3)$ & & & & & & &
\end{tabular}
\end{center}
\item For which values of $n$ do we get $1$ on both rows? What do they have in common?
\item State a conjecture. Test it for $n=13$ and $a=5$.
\item For $n=7$, compute $r_7(7)$ and $r_7(14)$. How should the conjecture be made more precise?
\end{enumerate}
\par\emph{Observing is not proving: the properties conjectured here are established in the course.}
\end{activity}
\begin{solution}[Activity]
\begin{enumerate}
\item $2^4=16\equiv1\pmod5$; $3^4=81\equiv1\pmod5$; $2^6=64\equiv1\pmod7$; $3^6=729=7\times104+1\equiv1\pmod7$.
\item Row $r_n(2)$: $1$, $2$, $1$, $0$, $4$, $2$, $1$. Row $r_n(3)$: $1$, $3$, $1$, $3$, $0$, $3$, $1$.
\item For $n=5$, $7$ and $11$: these are the prime numbers in the table.
\item Conjecture: if $p$ is prime, then $a^{p-1}\equiv1\pmod p$. For $p=13$ and $a=5$: $5^2=25\equiv-1$, so $5^{12}=\left(5^2\right)^6\equiv1\pmod{13}$.
\item $r_7(7)=r_7(14)=0$: we must assume that $a$ is not a multiple of $p$. This is Fermat's little theorem.
\end{enumerate}
\end{solution}
\begin{teachernote}[Running the activity]
Insist: the converse is false. The number $341=11\times31$ is not prime, and yet $2^{340}\equiv1\pmod{341}$ (have it checked with \texttt{pow}). Primality tests and Carmichael numbers are ``possible problems'' of the syllabus.
\end{teachernote}

\mxnewpage
\section{Course notes}
\subsection{Prime numbers}
\begin{definition}[Prime number]
A natural number $p$ is \emph{prime} if it has exactly two positive divisors, $1$ and $p$. Thus $1$ is not prime.
\end{definition}
\begin{theorem}[Smallest divisor]
Every integer $n\ge2$ has a prime divisor: its smallest divisor greater than or equal to $2$. If $n$ is not prime, this divisor $p$ satisfies $p^2\le n$.
\end{theorem}
\begin{proof}
Let $p$ be the smallest divisor of $n$ greater than or equal to $2$ (it exists, since $n$ is one). If $p$ were not prime, it would have a divisor $d$ with $2\le d<p$, which would also divide $n$: a contradiction. If $n$ is not prime, $n=pq$ with $q\ge2$ a divisor of $n$, so $q\ge p$ and $n=pq\ge p^2$.
\end{proof}
\begin{method}[Testing whether an integer is prime]
To know whether $n\ge2$ is prime, it is enough to check that no prime $p$ with $p^2\le n$ divides $n$. Example: $211$ is divisible neither by $2$, $3$, $5$, $7$, $11$, nor by $13$, and $17^2=289>211$: $211$ is prime.
\end{method}
\begin{theorem}[Infinitude of the primes]
There are infinitely many prime numbers.
\end{theorem}
\begin{proof}
Assume there are only finitely many: $p_1,\dots,p_k$. Let $N=p_1p_2\cdots p_k+1\ge2$. By the previous theorem, $N$ has a prime divisor, which is one of the $p_i$. This $p_i$ divides $N$ and the product $p_1\cdots p_k$, so it divides their difference $1$: this is impossible. The assumption is false.
\end{proof}
\begin{keypoint}[Sieve of Eratosthenes in Python]
Cross out the multiples of each prime starting from its square; the numbers not crossed out are prime.
\begin{lstlisting}
def crible(n):
    premier = [True] * (n + 1)
    premier[0] = premier[1] = False
    p = 2
    while p * p <= n:
        if premier[p]:
            for m in range(p * p, n + 1, p):
                premier[m] = False
        p = p + 1
    return [k for k in range(n + 1) if premier[k]]
\end{lstlisting}
There are $25$ prime numbers less than $100$.
\end{keypoint}
\begin{theorem}[Prime number and product]
Let $p$ be a prime number. If $p\mid ab$, then $p\mid a$ or $p\mid b$. More generally, if $p$ divides a product, it divides one of the factors.
\end{theorem}
\begin{proof}
If $p\nmid a$, then $\pgcd(p,a)=1$, because the only positive divisors of $p$ are $1$ and $p$. Gauss's theorem gives $p\mid b$. The case of a product of several factors follows by induction on their number.
\end{proof}

\subsection{Prime factorisation}
\begin{theorem}[Factorisation]
Every integer $n\ge2$ can be written as a product of primes, $n=p_1^{\alpha_1}p_2^{\alpha_2}\cdots p_k^{\alpha_k}$, with $p_1<p_2<\dots<p_k$ prime and $\alpha_i\ge1$. This expression is unique.
\end{theorem}
\begin{proof}
\emph{Existence}, by strong induction. The integer $2$ is prime. Assume every integer in $\ent{2}{n}$ can be factorised. If $n+1$ is prime, we are done; otherwise $n+1=ab$ with $2\le a,b\le n$, and we put the factorisations of $a$ and $b$ together.\par
\emph{Uniqueness.} If $p_1\cdots p_r=q_1\cdots q_s$ (repeated prime factors), $p_1$ divides the product of the $q_j$, hence one of them, which is prime: $p_1=q_j$. We cancel and start again; the two lists of factors coincide.
\end{proof}
\begin{keypoint}[Factorising in Python]
Divide by the smallest divisor $p\ge2$ as long as possible; it is prime. When $p^2>n$, the remaining quotient, if it exceeds $1$, is prime.
\begin{lstlisting}
def decomposition(n):
    facteurs = []
    p = 2
    while p * p <= n:
        while n % p == 0:
            facteurs.append(p)
            n = n // p
        p = p + 1
    if n > 1:
        facteurs.append(n)
    return facteurs
\end{lstlisting}
For example, \texttt{decomposition(360)} returns \texttt{[2, 2, 2, 3, 3, 5]}.
\end{keypoint}
\begin{theorem}[Divisors]
If $n=p_1^{\alpha_1}\cdots p_k^{\alpha_k}$, the positive divisors of $n$ are the integers $p_1^{\beta_1}\cdots p_k^{\beta_k}$ with $0\le\beta_i\le\alpha_i$. There are $(\alpha_1+1)(\alpha_2+1)\cdots(\alpha_k+1)$ of them.
\end{theorem}
\begin{method}[GCD by factorisation]
The gcd of two integers is the product of the common prime factors, each raised to the smaller of the two exponents. Example: $360=2^3\times3^2\times5$ and $84=2^2\times3\times7$ give $\pgcd(360,84)=2^2\times3=12$.
\end{method}

\subsection{Fermat's little theorem}
\begin{theorem}[Fermat's little theorem]
Let $p$ be a prime number and $a$ an integer. Then $a^p\equiv a\pmod p$. If moreover $p$ does not divide $a$, then $a^{p-1}\equiv1\pmod p$.
\end{theorem}
\begin{proof}
\emph{Step 1.} For $1\le k\le p-1$, $p$ divides $\binom pk$. Indeed $k\binom pk=p\binom{p-1}{k-1}$, so $p\mid k\binom pk$; since $1\le k<p$, $p$ does not divide $k$ and $\pgcd(p,k)=1$; Gauss gives $p\mid\binom pk$.\par
\emph{Step 2.} By the binomial formula, for every integer $a$, $(a+1)^p=a^p+\sum_{k=1}^{p-1}\binom pk a^k+1\equiv a^p+1\pmod p$.\par
\emph{Step 3.} Induction on $a\ge0$: $0^p\equiv0$, and if $a^p\equiv a$, then $(a+1)^p\equiv a^p+1\equiv a+1\pmod p$. For $a<0$, choose an integer $b\ge0$ such that $b\equiv a\pmod p$; then $a^p\equiv b^p\equiv b\equiv a$.\par
\emph{Step 4.} If $p\nmid a$, then $p\mid a^p-a=a\left(a^{p-1}-1\right)$ and $\pgcd(p,a)=1$; Gauss gives $p\mid a^{p-1}-1$.
\end{proof}
\begin{example}[A remainder]
Remainder of $3^{2026}$ in the division by $11$: $11$ is prime and does not divide $3$, so $3^{10}\equiv1\pmod{11}$. Since $\num{2026}=10\times202+6$, $3^{2026}\equiv3^6=729\equiv3\pmod{11}$ (because $729=66\times11+3$).
\end{example}
\begin{pitfall}
The converse is false: $2^{340}\equiv1\pmod{341}$, whereas $341=11\times31$.
\end{pitfall}

\mxnewpage
\section{Exercises}
\begin{exercise}[Prime or not?]
Are the following integers prime? $211$; $221$; $\num{1001}$; $257$.
\end{exercise}
\begin{solution}
$211$ is prime (course). $221=13\times17$. $\num{1001}=7\times11\times13$. For $257$: none of $2$, $3$, $5$, $7$, $11$, $13$ divides it and $17^2=289>257$: it is prime.
\end{solution}

\begin{exercise}[Factorisation and divisors]
Factorise $720$ into primes. How many positive divisors does it have? Are there many divisors of $720$ that are perfect squares?
\end{exercise}
\begin{solution}
$720=2^4\times3^2\times5$. Number of divisors: $(4+1)(2+1)(1+1)=30$. A square divisor has even exponents: $2^{\beta_1}3^{\beta_2}5^{\beta_3}$ with $\beta_1\in\ens{0;2;4}$, $\beta_2\in\ens{0;2}$, $\beta_3=0$; there are $6$ of them: $1$, $4$, $16$, $9$, $36$, $144$.
\end{solution}

\begin{beyond}[Least common multiple]The least positive common multiple of two nonzero integers, written lcm, is not part of the syllabus. It is used in the next exercise, which is not assessed.\end{beyond}
\begin{exercise}[GCD and LCM by factorisation]
Factorise $\num{1260}$ and $\num{1050}$, then deduce their gcd and their least common multiple. Check that the product of the gcd and the lcm equals the product of the two numbers.
\end{exercise}
\begin{solution}
$\num{1260}=2^2\times3^2\times5\times7$ and $\num{1050}=2\times3\times5^2\times7$. GCD: $2\times3\times5\times7=210$; LCM: $2^2\times3^2\times5^2\times7=\num{6300}$. And $210\times\num{6300}=\num{1323000}=\num{1260}\times\num{1050}$.
\end{solution}

\begin{exercise}[Fermat and remainders]
\begin{enumerate}
\item Determine the remainder of $2^{2026}$ in the division by $13$.
\item Prove that for every integer $n$, $n^{13}-n$ is divisible by $13$.
\end{enumerate}
\end{exercise}
\begin{solution}
\begin{enumerate}
\item $13$ is prime and does not divide $2$: $2^{12}\equiv1\pmod{13}$. $\num{2026}=12\times168+10$, so $2^{2026}\equiv2^{10}=\num{1024}\pmod{13}$. Now $\num{1024}=13\times78+10$: the remainder is $10$.
\item This is the first form of Fermat's little theorem, $n^{13}\equiv n\pmod{13}$.
\end{enumerate}
\end{solution}

\begin{exercise}[Divisible by $42$]
Prove that for every integer $n$, $n^7-n$ is divisible by $42$.
\end{exercise}
\begin{solution}
$42=2\times3\times7$. By Fermat, $n^7\equiv n\pmod7$. Modulo $2$: $n^2\equiv n$, so $n^7\equiv n$. Modulo $3$: $n^3\equiv n$, so $n^7=n^3\times n^3\times n\equiv n\times n\times n=n^3\equiv n$. Thus $2$, $3$ and $7$ divide $n^7-n$; they are pairwise coprime, so (corollary of Gauss's theorem, applied twice) $42\mid n^7-n$.
\end{solution}

\begin{exercise}[Square of a prime number]
Let $p$ be a prime number greater than or equal to $5$. Prove that $24$ divides $p^2-1$.
\end{exercise}
\begin{solution}
$p^2-1=(p-1)(p+1)$. Since $p$ is odd, $p-1$ and $p+1$ are two consecutive even integers: one of them is a multiple of $4$, so $8\mid p^2-1$. Since $p$ is not a multiple of $3$, $p\equiv1$ or $2\pmod3$, so $3\mid p-1$ or $3\mid p+1$. $8$ and $3$ are coprime: $24\mid p^2-1$.
\end{solution}

\begin{exercise}[Mersenne numbers]
\begin{enumerate}
\item Using the factorisation of $z^m-1$ (chapter 4), prove that if $n=ab$ with $a\ge2$ and $b\ge2$, then $2^a-1$ divides $2^n-1$.
\item Deduce that if $2^n-1$ is prime, then $n$ is prime.
\item Is the converse true? Compute $2^{11}-1$ and check that $23$ divides it.
\end{enumerate}
\end{exercise}
\begin{solution}
\begin{enumerate}
\item With $z=2^a$: $2^n-1=z^b-1=(z-1)\left(z^{b-1}+\dots+z+1\right)$, so $2^a-1\mid2^n-1$.
\item If $n\ge2$ is not prime, $n=ab$ as in 1, and $1<2^a-1<2^n-1$ is a divisor of $2^n-1$, which is therefore not prime. Finally $2^1-1=1$ is not prime. By contraposition, if $2^n-1$ is prime, $n$ is prime.
\item No: $11$ is prime, but $2^{11}-1=\num{2047}=23\times89$.
\end{enumerate}
\end{solution}

\mxnewpage
\section{Problem}
\begin{problem}[A miniature RSA cipher]
Choose two primes $p=5$ and $q=11$, and let $n=pq=55$ and $\varphi=(p-1)(q-1)=40$.
\begin{enumerate}
\item Choose $e=3$. Explain why $e$ and $\varphi$ are coprime, then find an integer $d\in\ent{1}{39}$ such that $ed\equiv1\pmod{40}$.
\item Encryption: a message is an integer $m\in\ent{0}{54}$; it is encrypted as the remainder $c$ of $m^e$ modulo $n$. Encrypt $m=8$.
\item Decryption: we take for granted for now that $c^d\equiv m\pmod n$. Check with a calculator or Python that $17^{27}\equiv8\pmod{55}$.
\item We prove that $m^{ed}\equiv m\pmod{55}$ for every integer $m$.
\begin{enumerate}
\item Write $ed=1+40k$ with $k$ a natural number.
\item If $5\nmid m$, prove that $m^{40k}\equiv1\pmod5$; deduce $m^{ed}\equiv m\pmod5$. What happens if $5\mid m$?
\item Establish likewise $m^{ed}\equiv m\pmod{11}$, then conclude.
\end{enumerate}
\item Why does the security of this cipher rest on the difficulty of factorising $n$ when $p$ and $q$ are very large?
\end{enumerate}
\end{problem}
\begin{solution}[Problem]
\begin{enumerate}
\item $40=2^3\times5$ and $3$ does not divide $40$: they are coprime. $3\times27=81=2\times40+1$: $d=27$.
\item $8^3=512=9\times55+17$: $c=17$.
\item We do get $8$ (for example \texttt{pow(17, 27, 55)}).
\item
\begin{enumerate}
\item $ed=81=1+40\times2$: $k=2$.
\item By Fermat, $m^4\equiv1\pmod5$, so $m^{40k}=\left(m^4\right)^{10k}\equiv1$ and $m^{ed}=m\times m^{40k}\equiv m\pmod5$. If $5\mid m$, both sides are congruent to $0$.
\item Modulo $11$: if $11\nmid m$, $m^{10}\equiv1$, so $m^{40k}=\left(m^{10}\right)^{4k}\equiv1$; if $11\mid m$, both sides are zero modulo $11$. Thus $5$ and $11$ divide $m^{ed}-m$ and are coprime: $55\mid m^{ed}-m$.
\end{enumerate}
\item Computing $d$ by this method requires $\varphi=(p-1)(q-1)$, hence $p$ and $q$. Knowing only $n$ and $e$, factorising $n$ would be enough to find $d$. For well-chosen primes with several hundred digits, no known method factorises $n$ in a reasonable time on current computers. The security of RSA rests on this assumed difficulty: it has not been proved that every way of decrypting goes through factorising $n$.
\end{enumerate}
\end{solution}
\begin{teachernote}[Status of the problem]
RSA is a ``possible problem'' of the syllabus. Have the encryption and decryption functions written in Python and tested on the $55$ possible messages. Question 5 gives an idea of security; it does not call for a proof.
\end{teachernote}
